% ----------------------------------------------------------------------------
%  3. Problem Statement
% ----------------------------------------------------------------------------
\section{Problem Statement}
\label{sec:problem}

\subsection{What we have}

\begin{itemize}
  \item A list of \textbf{patients}, each with an acuity (CRIT,
        HIGH, MED, LOW), an age, and a diagnosis.
  \item A list of \textbf{workers}, each with a role (PGY1, INTERN,
        NURSE) and a finite shift in minutes.
  \item Three reference catalogues:
        \begin{itemize}
          \item 418 diseases, each mapping to a list of required
                tests and a base priority.
          \item 183 tests, each with a required role and a
                duration in minutes.
          \item 325 ``A must precede B'' edges across the test
                catalogue.
        \end{itemize}
\end{itemize}

\subsection{What ``a schedule'' means}

A schedule says, for every (patient, test) pair: who does it, when
they start. To be valid, the schedule must:

\begin{itemize}
  \item Use a worker whose role matches the test.
  \item Respect every ``must precede'' edge.
  \item Never double-book a worker.
  \item Cover every task whose required role exists on the roster.
\end{itemize}

\subsection{What ``a good schedule'' means}

Three smaller-is-better scores, all reported separately:

\begin{itemize}
  \item Average finish time of urgent (CRIT+HIGH) patients.
  \item Total schedule length.
  \item Worst single-worker overtime.
\end{itemize}

\noindent
A schedule's ``feasibility verdict'' is just: total overtime
$\leq$ a threshold the operator sets.

\subsection{What we don't try to handle}

\begin{itemize}
  \item Patient deteriorating mid-shift (acuity is fixed once set).
  \item Random or unpredictable test durations (for example, if
        the cardiology ward sends a patient for an Echocardiogram,
        the wait time depends on the inflow from other wards
        requesting the same Echocardiogram service).
  \item Worker breaks or training time.
  \item Per-(patient, test) priority overrides.
\end{itemize}
